\uuid{pNqE}
\exo7id{4609}
\titre{exo7 4609}
\auteur{quercia}
\organisation{exo7}
\datecreate{2010-03-14}
\isIndication{false}
\isCorrection{true}
\chapitre{Série entière}
\sousChapitre{Intégrales}
\cpge{mpsi-mp2i.fonctions-valeurs-reelles-complexes.c, pcsi.techniques-calcul-fonctions.c, ptsi.techniques-calcul-fonctions.c, mpsi-mp2i.integration.c, pcsi.integration.b, ptsi.integration.b, bcpst1.fonctions-reelles-usuelles.c, bcpst1.derivees-primitives-integrales.d, mp-mpi.series-entieres.a, pc.series-entieres.a, psi.series-entieres.a, pt.series-entieres.a, mp-mpi.series-entieres.d, pc.series-entieres.c, psi.series-entieres.c, pt.series-entieres.c, mp-mpi.suites-series-fonctions.e, pc.suites-series-fonctions.c, psi.suites-series-fonctions.c, bcpst2.integrales-generalisees}
\module{Analyse}
\niveau{L2}
\difficulte{}

\contenu{
\question{
Établir la convergence puis calculer la valeur de
$ \int_{t=0}^1 \frac{\ln(t^2)\ln(1-t^2)}{t^2}\,d t$.
}
\reponse{
Développer en série entière $\ln(1-t^2)$.
$I = \frac{\pi^2}2 - 4\ln 2$.
}
}
